\ifdefined\gkver\else\def\gkver{student}\fi
\documentclass[\gkver]{gaokaozhenti}
\begin{document}

\chapter{2014年安徽理综}
%% paperType: 理综物理部分

\begin{enumerate}
\item
%% number: 14
%% paperName: 2014年安徽理综第 14 题 6 分
%% typeId: 1
%% type: 单选题
%% chapter: 机械振动
%% point: 单摆
%% method: 无
%% score: 6
%% degree: 450
%% duplicateId: 0
%% body:
在科学研究中，科学家常将未知现象同已知现象进行比较，找出其共同点，进一步推测未知现象的特性和规律。法国物理学家库仑在研究异种电荷的吸引问题时，曾将扭秤的振动周期与电荷间距离的关系类比单摆的振动周期与摆球到地心距离的关系。已知单摆摆长为 $l$，引力常量为 $G$。地球的质量为 $M$。摆球到地心的距离为 $r$，则单摆振动周期 $T$ 与距离 $r$ 的关系式为\xzanswer{B}
\fourchoices[answer=B]
{$T=2\pi r\sqrt{\frac{GM}{l}}$}
{$T=2\pi r\sqrt{\frac{l}{GM}}$}
{$T=\frac{2\pi}{r}\sqrt{\frac{GM}{l}}$}
{$T=2\pi l\sqrt{\frac{l}{GM}}$}
%% answer: B
%% memo:
\memoanswer{
由于万有引力使物体产生加速度，由牛顿第二定律得：$G\frac{Mm}{r^{2}}=mg$，而单摆的振动周期公式为 $T=2\pi\sqrt{\frac{l}{g}}$，联立得：$T=2\pi r\sqrt{\frac{l}{GM}}$。B 正确。
}

\item
%% number: 15
%% paperName: 2014年安徽理综第 15 题 6 分
%% typeId: 1
%% type: 单选题
%% chapter: 功和能
%% point: 机械能守恒定律
%% method: 图象法
%% score: 6
%% degree: 450
%% duplicateId: 0
%% body:
如图所示，有一内壁光滑的闭合椭圆形管道，置于竖直平面内，MN 是通过椭圆中心 O 点的水平线。已知一小球从 M 点出发，初速率为 $v_{0}$，沿管道 MPN 运动，到 N 点的速率为 $v_{1}$，所需的时间为 $t_{1}$；若该小球仍由 M 点以初速率 $v_{0}$ 出发，而沿管道 MQN 运动，到 N 点的速率为 $v_{2}$，所需时间为 $t_{2}$。则\xzanswer{A}
\onepicture[w=4.2cm, align=right]{figs/15.png}
\fourchoices[answer=A]
{$v_{1}=v_{2}$，$t_{1}>t_{2}$}
{$v_{1}<v_{2}$，$t_{1}>t_{2}$}
{$v_{1}=v_{2}$，$t_{1}<t_{2}$}
{$v_{1}<v_{2}$，$t_{1}<t_{2}$}
%% answer: A
%% memo:
\memoanswer{
由于是内壁光滑的闭合椭圆形管道，运动中只有重力做功，机械能守恒，MON 在同一水平线上，故 $v_{1}=v_{2}=v_{0}$；而沿管道 MPN 运动，先减速后加速，沿管道 MQN 运动，先加速后减速，前者平均速率小，后者平均速率大，运动的路程相同，故 $t_{1}>t_{2}$。A 正确。
}

\item
%% number: 16
%% paperName: 2014年安徽理综第 16 题 6 分
%% typeId: 1
%% type: 单选题
%% chapter: 机械振动
%% point: 振动图象
%% method: 无
%% score: 6
%% degree: 450
%% duplicateId: 0
%% body:
一简谐横波沿 $x$ 轴正向传播，图 1 是 $t=0$ 时刻的波形图，图 2 是介质中某点的振动图象，则该质点的 $x$ 坐标值合理的是\xzanswer{C}
\onepicture[w=9cm]{figs/16.png}
\fourchoices[answer=C]
{$0.5\Um$}
{$1.5\Um$}
{$2.5\Um$}
{$3.5\Um$}
%% answer: C
%% memo:
\memoanswer{
由图 2 结合图 1 可知该质点 $x$ 坐标值可能是 $1.5\Um$ 和 $2.5\Um$，而简谐横波沿 $x$ 轴正向传播，由图 1 可得向下振动的质点为 $x$ 坐标值 $2.5\Um$ 的质点，故 C 正确。
}

\item
%% number: 17
%% paperName: 2014年安徽理综第 17 题 6 分
%% typeId: 1
%% type: 单选题
%% chapter: 电场
%% point: 电场中的图象
%% method: 无
%% score: 6
%% degree: 450
%% duplicateId: 0
%% body:
一带电粒子在电场中仅受静电力作用，做初速度为零的直线运动。取该直线为 $x$ 轴，起始点 O 为坐标原点，其电势能 $E_{P}$ 与位移 $x$ 的关系如右图所示。下列图象中合理的是\xzanswer{D}
\onepicture[w=4.6cm, align=right]{figs/17a.png}
\fourchoices[answer=D, ispicture=true, h=2.0cm]
{figs/17b.png}
{figs/17c.png}
{figs/17d.png}
{figs/17e.png}
%% answer: D
%% memo:
\memoanswer{
由电场力做功与电势能的关系：$Fx=-\Delta E_{p}$，可知 $E_{P}-x$ 图线的斜率表示静电力 $F$ 的大小，可见静电力 $F$ 逐渐减小，而 $F=qE$，故不是匀强电场，A 错误；根据牛顿第二定律粒子做加速度减小的加速运动，C 错误，D 正确；根据能量守恒 $\Delta E_{k}=-\Delta E_{p}$，比较图线 B 错误。正确选项 D。
}

\item
%% number: 18
%% paperName: 2014年安徽理综第 18 题 6 分
%% typeId: 1
%% type: 单选题
%% chapter: 磁场
%% point: 带电粒子在磁场中的运动
%% method: 无
%% score: 6
%% degree: 450
%% duplicateId: 0
%% body:
“人造小太阳”托卡马克装置使用强磁场约束高温等离子体，使其中的带电粒子被尽可能限制在装置内部，而不与装置器壁碰撞。已知等离子体中带电粒子的平均动能与等离子体的温度 $T$ 成正比，为约束更高温度的等离子体，则需要更强的磁场，以使带电粒子的运动半径不变。由此可判断所需的磁感应强度 $B$ 正比于\xzanswer{A}
\fourchoices[answer=A]
{$\sqrt T$}
{$T$}
{$\sqrt{T^{3}}$}
{$T^{2}$}
%% answer: A
%% memo:
\memoanswer{
由于等离子体中带电粒子的平均动能与等离子体的温度 $T$ 成正比，即 $\overline{E_{k}}\propto T$。带电粒子在磁场中做圆周运动，洛伦兹力提供向心力：$qvB=m\frac{v^{2}}{R}$，得 $B=\frac{mv}{qR}$。而 $E_{k}=\frac{1}{2}mv^{2}$，故可得：$B=\frac{\sqrt{2mE_{k}}}{qR}$。又带电粒子的运动半径不变，所以 $B\propto\sqrt{E_{k}}\propto\sqrt T$。A 正确。
}

\item
%% number: 19
%% paperName: 2014年安徽理综第 19 题 6 分
%% typeId: 1
%% type: 单选题
%% chapter: 曲线运动
%% point: 圆周运动的应用
%% method: 无
%% score: 6
%% degree: 450
%% duplicateId: 0
%% body:
如图所示，一倾斜的匀质圆盘绕垂直于盘面的固定对称轴以恒定的角速度 $\omega$ 转动，盘面上离转轴距离 $2.5\Um$ 处有一小物体与圆盘始终保持相对静止。物体与盘面间的动摩擦因数为 $\frac{\sqrt 3}{2}$（设最大静摩擦力等于滑动摩擦力），盘面与水平面的夹角为 $30^\circ$，$g$ 取 $10\Umsq$。则 $\omega$ 的最大值是\xzanswer{C}
\onepicture[w=4.2cm, align=right]{figs/19.png}
\fourchoices[answer=C]
{$\sqrt 5\Urads$}
{$\sqrt 3\Urads$}
{$1.0\Urads$}
{$0.5\Urads$}
%% answer: C
%% memo:
\memoanswer{
由于小物体随匀质圆盘做圆周运动，其向心力由小物体受到的指向圆心的合力提供，在最下端时指向圆心的合力最小。根据牛顿第二定律：$f-mg\sin 30^\circ=m\omega^{2}r$，又 $f\leq f_{m}=\mu mg\cos 30^\circ$，解得 $\omega\leq 1.0\Urads$，要使小物体与圆盘始终保持相对静止，则 $\omega$ 的最大值是 $1.0\Urads$。C 正确。
}

\item
%% number: 20
%% paperName: 2014年安徽理综第 20 题 6 分
%% typeId: 1
%% type: 单选题
%% chapter: 电磁感应
%% point: 电磁感应能量分析
%% method: 无
%% score: 6
%% degree: 450
%% duplicateId: 0
%% body:
英国物理学家麦克斯韦认为，磁场变化时会在空间激发感生电场。如图所示，一个半径为 $r$ 的绝缘细圆环水平放置，环内存在竖直向上的匀强磁场 $B$，环上套一带电量为 $+q$ 的小球。已知磁感应强度 $B$ 随时间均匀增加，其变化率为 $k$，若小球在环上运动一周，则感生电场对小球的作用力所做功的大小是\xzanswer{D}
\onepicture[w=4.4cm, align=right]{figs/20.png}
\fourchoices[answer=D]
{0}
{$\frac{1}{2}r^{2}qk$}
{$2\pi r^{2}qk$}
{$\pi r^{2}qk$}
%% answer: D
%% memo:
\memoanswer{
由法拉第电磁感应定律得感生电动势：$E=\frac{\Delta \Phi}{\Delta t}=\frac{\Delta B}{\Delta t}\pi r^{2}=k\pi r^{2}$，而电场力做功 $W=qU$，小球在环上运动一周 $U=E$，故 $W=\pi r^{2}qk$。D 正确。
}

\item
%% number: 21
%% paperName: 2014年安徽理综第 21 题 10 分
%% typeId: 4
%% type: 实验题
%% chapter: 电路
%% point: 电路实验
%% method: 无
%% score: 10
%% degree: 450
%% duplicateId: 0
%% body:
\begin{enumerate}
	\item
	Ⅰ．图 1 是“研究平抛物体运动”的实验装置图，通过描点画出平抛小球的运动轨迹。

	（1）以下是实验过程中的一些做法，其中合理的有\tkanswer[1.5cm]{ac}。

	（A）安装斜槽轨道，使其末端保持水平

	（B）每次小球释放的初始位置可以任意选择

	（C）每次小球应从同一高度由静止释放

	（D）为描出小球的运动轨迹，描绘的点可以用折线连接

	（2）实验得到平抛小球的运动轨迹，在轨迹上取一些点，以平抛起点 $O$ 为坐标原点，测量它们的水平坐标 $x$ 和竖直坐标 $y$，图 2 中 $y-x^{2}$ 图象能说明平抛小球运动轨迹为抛物线的是\tkanswer[1cm]{c}。

	（3）图 3 是某同学根据实验画出的平抛小球的运动轨迹，$O$ 为平抛的起点，在轨迹上任取三点 $A$、$B$、$C$，测得 $A$、$B$ 两点竖直坐标 $y_{1}$ 为 $5.0\Ucm$、$y_{2}$ 为 $45.0\Ucm$，$A$、$B$ 两点水平间距 $\Delta x$ 为 $40.0\Ucm$。则平抛小球的初速度 $v_{0}$ 为\tkanswer[1.5cm]{2.0}$\Ums$，若 $C$ 点的竖直坐标 $y_{3}$ 为 $60.0\Ucm$，则小球在 $C$ 点的速度 $v_{C}$ 为\tkanswer[1.5cm]{4.0}$\Ums$（结果保留两位有效数字，$g$ 取 $10\Umsq$）。

	\threepicture[labelA=2014安徽理综21a, labelB=2014安徽理综21b, labelC=2014安徽理综21c, numA=1, numB=2, numC=3, showlabel=false, wA=3.8cm, wB=7.0cm, wC=3.2cm, align=right]
	{figs/21a.png}
	{figs/21b.png}
	{figs/21c.png}

	\item
	Ⅱ．某同学为了测量一个量程为 $3\UV$ 的电压表的内阻，进行了如下实验：

	（1）他先用多用电表进行了正确的测量，测量时指针位置如图 1 所示，得出电压表内阻为 $3.00\times10^{3}\UO$，此时电压表的指针也偏转了。已知多用表欧姆档表盘中央刻度值为“15”，表内电池电动势为 $1.5\UV$，则电压表的示数应为\tkanswer[1.2cm]{1.0}$\UV$（结果保留两位有效数字）。

	（2）为了更准确地测量该电压表的内阻 $R_{V}$，该同学设计了图 2 所示的电路图，实验步骤如下：

	A．断开开关 $S$，按图 2 连接好电路；

	B．把滑动变阻器 $R$ 的滑片 $P$ 滑到 $b$ 端；

	C．将电阻箱 $R_{0}$ 的阻值调到零；

	D．闭合开关 $S$；

	E．移动滑动变阻器 $R$ 的滑片 $P$ 的位置，使电压表的指针指到 $3\UV$ 位置；

	F．保持滑动变阻器 $R$ 的滑片 $P$ 位置不变，调节电阻箱 $R_{0}$ 的阻值使电压表的指针指到 $1.5\UV$ 位置，读出此时电阻箱 $R_{0}$ 的阻值，此值即为电压表内阻 $R_{V}$ 的测量值；

	G．断开开关 $S$。

	实验中可供选择的实验器材有：

	a．待测电压表

	b．滑动变阻器：最大阻值 $2000\UO$

	c．滑动变阻器：最大阻值 $10\UO$

	d．电阻箱：最大阻值 $9999.9\UO$，阻值最小改变量为 $0.1\UO$

	e．电阻箱：最大阻值 $999.9\UO$，阻值最小改变量为 $0.1\UO$

	f．电池组：电动势约 $6\UV$，内阻可忽略

	g．开关，导线若干

	按照这位同学设计的实验方案，回答下列问题：

	①要使测量更精确，除了选用电池组、导线、开关和待测电压表外，还应从提供的滑动变阻器中选用\tkanswer[1cm]{c}（填“b”或“c”），电阻箱中选用\tkanswer[1cm]{d}（填“d”或“e”）。

	②电压表的内阻 $R_{V}$ 的测量值 $R_{\text{测}}$ 和真实值 $R_{\text{真}}$ 相比，$R_{\text{测}}$\tkanswer[0.8cm]{$>$}$R_{\text{真}}$（填“$>$”或“$<$”）；若 $R_{V}$ 越大，则 $\frac{|R_{\text{测}}-R_{\text{真}}|}{R_{\text{真}}}$ 越\tkanswer[1.2cm]{小}（填“大”或“小”）。

	\twopicture[labelA=2014安徽理综21d, labelB=2014安徽理综21e, numA=1, numB=2, wA=4.0cm, wB=4.0cm, align=right]
	{figs/21d.png}
	{figs/21e.png}
\end{enumerate}
%% answer: 
\jdanswer{
\begin{enumerate}
	\item
	Ⅰ．（1）ac；（2）c；（3）$2.0\Ums$，$4.0\Ums$。

	\item
	Ⅱ．（1）$1.0\UV$；（2）①c，d；②$>$，小。

\end{enumerate}
}
%% memo:
\memoanswer{
（1）“研究平抛物体运动”的实验斜槽轨道末端保持水平为了保证水平初速度。从同一高度由静止释放为了保证每次使用水平初速度相同。a、c 正确。

（2）平抛物体运动规律：$x=v_{0}t$，$y=\frac{1}{2}gt^{2}$，得 $y=\frac{g}{2v_{0}^{2}}x^{2}$，$y-x^{2}$ 图象是一条倾斜直线。c 正确。

（3）由 $y=\frac{1}{2}gt^{2}$ 得 $t=\sqrt{\frac{2y}{g}}$，则 $t_{1}=0.1\Us$、$t_{2}=0.3\Us$，所以平抛小球的初速度 $v_{0}=\frac{\Delta x}{t_{2}-t_{1}}=2.0\Ums$。而 $v_{Cy}=\sqrt{2gy_{3}}=2\sqrt 3\Ums$，故 $C$ 点的速度 $v_{C}=\sqrt{v_{0}^{2}+v_{Cy}^{2}}=4.0\Ums$。

（1）多用电表的中值电阻即为其内阻，内阻为 $1.50\times10^{3}\UO$，与内阻为 $3.00\times10^{3}\UO$ 的电压表串联，电源电动势为 $1.5\UV$，可得电压表的读数即为 $R_{V}$ 两端电压为 $1.0\UV$。

（2）①由于滑动变阻器的分压影响产生误差，故选择阻值小的 c 能减小误差；又由于使用半偏法测电压表内阻，电阻箱最大阻值应大于 $3.00\times10^{3}\UO$，故选择 d。

②$R_{0}=0$ 时，电压表的指针指到 $3.0\UV$ 位置，电压表中的电流 $I_{1}=\frac{U_{1}}{R_{V}}$（$U_{1}=3.0\UV$）；使电压表的指针指到 $1.5\UV$ 位置时，电阻箱的电阻为 $R_{0}$，电压表中的电流 $I_{2}=\frac{U_{2}}{R_{V}+R_{0}}=\frac{U_{1}/2}{R_{V}}$，而滑动变阻器的分压影响 $U_{2}<U_{1}$，故 $R_{0}<R_{V}$ 即 $R_{\text{测}}<R_{\text{真}}$；$R_{V}$ 越大，滑动变阻器的分压影响越小，$R_{\text{测}}$ 越接近 $R_{\text{真}}$，故 $\frac{|R_{\text{测}}-R_{\text{真}}|}{R_{\text{真}}}$ 越小。
}

\item
%% number: 22
%% paperName: 2014年安徽理综第 22 题 14 分
%% typeId: 6
%% type: 计算题
%% chapter: 电场
%% point: 带电粒子的直线运动
%% method: 无
%% score: 14
%% degree: 450
%% duplicateId: 0
%% body:
如图所示，充电后的平行板电容器水平放置，电容为 $C$，极板间距离为 $d$，上极板正中有一小孔。质量为 $m$、电荷量为 $+q$ 的小球从小孔正上方高 $h$ 处由静止开始下落，穿过小孔到达下极板处速度恰为零（空气阻力忽略不计，极板间电场可视为匀强电场，重力加速度为 $g$）。求：
\begin{enumerate}
	\item
	小球到达小孔处的速度；
	\item
	极板间电场强度大小和电容器所带电荷量；
	\item
	小球从开始下落运动到下极板的时间。
\end{enumerate}
\onepicture[align=right, w=3.6cm]{figs/22.png}
%% answer: 
\jdanswer{
\begin{enumerate}
	\item
	$\sqrt{2gh}$
	\item
	$\frac{mg(h+d)}{qd}$，$C\frac{mg(h+d)}{q}$
	\item
	$\frac{h+d}{h}\sqrt{\frac{2h}{g}}$
\end{enumerate}
}
%% memo:
\memoanswer{
（1）由 $v^{2}=2gh$ 得 $v=\sqrt{2gh}$。

（2）在极板间带电小球受重力和电场力，有 $mg-qE=ma$，$0-v^{2}=2ad$，得 $E=\frac{mg(h+d)}{qd}$；由 $U=Ed$，$Q=CU$，得 $Q=C\frac{mg(h+d)}{q}$。

（3）由 $h=\frac{1}{2}gt_{1}^{2}$，$0=v+at_{2}$，综合可得 $t=\frac{h+d}{h}\sqrt{\frac{2h}{g}}$。
}

\item
%% number: 23
%% paperName: 2014年安徽理综第 23 题 16 分
%% typeId: 6
%% type: 计算题
%% chapter: 电磁感应
%% point: 电磁感应能量分析
%% method: 无
%% score: 16
%% degree: 450
%% duplicateId: 0
%% body:
如图 1 所示，匀强磁场的磁感应强度 $B$ 为 $0.5\UT$，其方向垂直于倾角 $\theta$ 为 $30^\circ$ 的斜面向上。绝缘斜面上固定有“$\Lambda$”形状的光滑金属导轨 MPN（电阻忽略不计），MP 和 NP 长度均为 $2.5\Um$。MN 连线水平，长为 $3\Um$。以 MN 的中点 O 为原点、OP 为 $x$ 轴建立一坐标系 $Ox$。一根粗细均匀的金属杆 CD，长度 $d$ 为 $3\Um$，质量 $m$ 为 $1\Ukg$，电阻 $R$ 为 $0.3\UO$，在拉力 $F$ 的作用下，从 MN 处以恒定的速度 $v=1\Ums$ 在导轨上沿 $x$ 轴正向运动（金属杆与导轨接触良好）。$g$ 取 $10\Umsq$。
\begin{enumerate}
	\item
	求金属杆 CD 运动过程中产生的感应电动势 $E$ 及运动到 $x=0.8\Um$ 电势差 $U_{CD}$；
	\item
	推导金属杆 CD 从 MN 处运动到 P 点过程中拉力 $F$ 与位置坐标 $x$ 的关系式，并在图 2 中画出 $F-x$ 关系图象；
	\item
	求金属杆 CD 从 MN 处运动到 P 点的全过程产生的焦耳热。
\end{enumerate}
\onepicture[align=right, w=13cm]{figs/23.png}
%% answer: 
\jdanswer{
\begin{enumerate}
	\item
	$1.5\UV$，$-0.6\UV$
	\item
	$F=12.5-3.75x$（$0\leq x\leq 2$），图象见解析
	\item
	$7.5\UJ$
\end{enumerate}
}
%% memo:
\memoanswer{
（1）金属杆 CD 在匀速运动中产生的感应电动势 $E=Blv$（$l=d$），$E=1.5\UV$（D 点电势高）。

当 $x=0.8\Um$ 时，金属杆在导轨间的电势差为零。设此时杆在导轨外的长度为 $l_{\text{外}}$，则 $l_{\text{外}}=d-\frac{OP-x}{OP}d$，$OP=\sqrt{MP^{2}-\left(\frac{MN}{2}\right)^{2}}$，得 $l_{\text{外}}=1.2\Um$。

由楞次定律判断 D 点电势高，故 CD 两端电势差 $U_{CD}=-Bl_{\text{外}}v$，$U_{CD}=-0.6\UV$。

（2）杆在导轨间的长度 $l$ 与位置 $x$ 的关系是 $l=\frac{OP-x}{OP}d=3-\frac{3}{2}x$，对应的电阻 $R_{l}$ 为 $R_{l}=\frac{l}{d}R$，电流 $I=\frac{Blv}{R_{l}}$，杆受安培力 $F_{\text{安}}$ 为 $F_{\text{安}}=BIl=7.5-3.75x$。根据平衡条件得 $F=F_{\text{安}}+mg\sin\theta$，即 $F=12.5-3.75x$（$0\leq x\leq 2$）。

画出的 $F-x$ 图象如图所示。

\onepicture[w=6cm]{figs/23_an.png}

（3）外力 $F$ 所做的功 $W_{F}$ 等于 $F-x$ 图线下所围成的面积，即 $W_{F}=\frac{5+12.5}{2}\times2\UJ=17.5\UJ$。而杆的重力势能增加量 $\Delta E_{P}=mg\overline{OP}\sin\theta$，故全过程产生的焦耳热 $Q=W_{F}-\Delta E_{P}=7.5\UJ$。
}

\item
%% number: 24
%% paperName: 2014年安徽理综第 24 题 20 分
%% typeId: 6
%% type: 计算题
%% chapter: 运动和力
%% point: 动量守恒定律
%% method: 无
%% score: 20
%% degree: 250
%% duplicateId: 0
%% body:
光滑水平地面上有一凹槽 $A$，中央放一小物块 $B$。物块与左右两边槽壁的距离如图所示，$L$ 为 $1.0\Um$。凹槽与物块的质量均为 $m$，两者之间的动摩擦因数 $\mu$ 为 $0.5$。开始时物块静止，凹槽以 $v_{0}=5\Ums$ 初速度向右运动，设物块与凹槽壁碰撞过程中没有能量损失，且碰撞时间不计。$g$ 取 $10\Umsq$。求：
\begin{enumerate}
	\item
	物块与凹槽相对静止时的共同速度；
	\item
	从凹槽开始运动到两者相对静止物块与右侧槽壁碰撞的次数；
	\item
	从凹槽开始运动到两者相对静止所经历的时间及该时间内凹槽运动的位移大小。
\end{enumerate}
\onepicture[align=right, w=8cm]{figs/24.png}
%% answer: 
\jdanswer{
\begin{enumerate}
	\item
	$2.5\Ums$
	\item
	6
	\item
	$5\Us$，$12.75\Um$
\end{enumerate}
}
%% memo:
\memoanswer{
（1）设两者间相对静止时的速度为 $v$，由动量守恒定律得 $mv_{0}=2mv$，$v=2.5\Ums$。

（2）物块与凹槽间的滑动摩擦力 $f=\mu N=\mu mg$。设两者间相对静止时的路程为 $s_{1}$，由动能定理得 $-fs_{1}=\frac{1}{2}(m+m)v^{2}-\frac{1}{2}mv_{0}^{2}$，得 $s_{1}=12.5\Um$。已知 $L=1\Um$，可推知物块与右侧槽壁共发生 6 次碰撞。

（3）设凹槽与物块碰前的速度分别为 $v_{1}$、$v_{2}$，碰后的速度分别为 $v_{1}^{\prime}$、$v_{2}^{\prime}$。有 $mv_{1}+mv_{2}=mv_{1}^{\prime}+mv_{2}^{\prime}$，$\frac{1}{2}mv_{1}^{2}+\frac{1}{2}mv_{2}^{2}=\frac{1}{2}mv_{1}^{\prime2}+\frac{1}{2}mv_{2}^{\prime2}$，得 $v_{1}^{\prime}=v_{2}$，$v_{2}^{\prime}=v_{1}$，即每碰撞一次凹槽与物块发生一次速度交换，在同一坐标系上两者的速度图线如图所示，根据碰撞次数可分为 13 段，凹槽、物块的 $v-t$ 图象在两条连续的匀变速运动图线间转换，故可用匀变速直线运动规律求时间。则 $v=v_{0}-at$，$a=\mu g$，解得 $t=5\Us$。

\onepicture[w=8cm]{figs/24_an.png}

凹槽的 $v-t$ 图象所包围的阴影面积即为凹槽的位移大小 $s_{2}$（等腰三角形面积共分 13 份，第一份面积为 $0.5L$，其余每份面积均为 $L$）。$s_{2}=\frac{1}{2}\frac{v_{0}}{2}t+6.5L$，$s_{2}=12.75\Um$。
}

\end{enumerate}

\end{document}
